% GENERATED FILE. Edit canonical lessons, then run npm run generate:books.
% canonical-source-hash: fnv1a64:ec97558105fdd665
% canonical-lessons: LA-C518-latus, LA-C518-coxa, LA-C518-calx, LA-C518-planta, LA-C518-umbilicus

\chapter{Side, hip and foot}
\label{ch:la-side-hip-and-foot}
\hlchaptermodality{518}

\begin{chapteropening}
\textbf{By the end of this chapter:} \emph{I can name the side, the hip, the heel and the sole.}

Complete the last lesson of chapter 518: I can name the side, the hip, the heel and the sole.
\end{chapteropening}

\section[latus, lateris]{latus, lateris — side, flank}

\label{lesson:LA-C518-latus}



\begin{quote}
Before the new one: say the Latin for sweat, then the Latin for a beat.
\end{quote}

\subsection*{You'll want to know: latus, lateris}
\textbf{latus, lateris} — \textquotedblleft{}side, flank\textquotedblright{}.

\begin{quote}
Dolor in latere sinistrō est.
\end{quote}

\begin{quote}
\emph{The pain is in the left side.}
\end{quote}

Neuter, third declension. Not the adjective \emph{lātus} — wide. English \emph{lateral}.

\begin{cousinweb}
\textbf{latus} is a native Latin noun; Lewis \& Short compares Greek \emph{platus} \textquotedblleft{}broad\textquotedblright{}, but the origin is not settled. English \emph{lateral} comes from it.
\end{cousinweb}

\begin{grammarlens}[title={What you've built}]
One more word for telling, explaining and describing at B1.
\end{grammarlens}

\subsection*{Your turn}
\begin{itemize}
\raggedright
  \item \emph{Say it:} \emph{latus, lateris}
  \item \emph{Say it:} \emph{latus, lateris}, once more
  \item \emph{Say it:} a sentence of your own with \emph{latus, lateris}
  \item \emph{Recall:} say the Latin for sweat, then the Latin for a beat, then say \emph{latus, lateris} again
\end{itemize}

\begin{tcolorbox}[breakable,colback=teal!4,colframe=teal!35!black,title={Before you move on}]
What does latus, lateris mean? (\textquotedblleft{}Side, flank\textquotedblright{}.) Say it once more.
\end{tcolorbox}
\section[coxa, coxae]{coxa, coxae — hip}

\label{lesson:LA-C518-coxa}



\begin{quote}
Before the new one: say the Latin for a beat, then the Latin for side.
\end{quote}

\subsection*{You'll want to know: coxa, coxae}
\textbf{coxa, coxae} — \textquotedblleft{}hip\textquotedblright{}.

\begin{quote}
Senex in coxā dolōrem sentit.
\end{quote}

\begin{quote}
\emph{The old man feels pain in his hip.}
\end{quote}

Feminine, first declension. French \emph{cuisse} and Italian \emph{coscia} — thigh — come from it.

\begin{cousinweb}
\textbf{coxa} is inherited from a Proto-Indo-European word for \textquotedblleft{}hip, joint\textquotedblright{}, and Old Irish \emph{cos} \textquotedblleft{}foot, leg\textquotedblright{} is a cousin (EtymDB). French \emph{cuisse} \textquotedblleft{}thigh\textquotedblright{} comes from it.
\end{cousinweb}

\begin{grammarlens}[title={What you've built}]
One more word for telling, explaining and describing at B1.
\end{grammarlens}

\subsection*{Your turn}
\begin{itemize}
\raggedright
  \item \emph{Say it:} \emph{coxa, coxae}
  \item \emph{Say it:} \emph{coxa, coxae}, once more
  \item \emph{Say it:} a sentence of your own with \emph{coxa, coxae}
  \item \emph{Recall:} say the Latin for a beat, then the Latin for side, then say \emph{coxa, coxae} again
\end{itemize}

\begin{tcolorbox}[breakable,colback=teal!4,colframe=teal!35!black,title={Before you move on}]
What does coxa, coxae mean? (\textquotedblleft{}Hip\textquotedblright{}.) Say it once more.
\end{tcolorbox}
\section[calx, calcis]{calx, calcis — heel}

\label{lesson:LA-C518-calx}



\begin{quote}
Before the new one: say the Latin for side, then the Latin for hip.
\end{quote}

\subsection*{You'll want to know: calx, calcis}
\textbf{calx, calcis} — \textquotedblleft{}heel\textquotedblright{}.

\begin{quote}
Calx mea post longum iter dolet.
\end{quote}

\begin{quote}
\emph{My heel hurts after the long journey.}
\end{quote}

Feminine, third declension. Another \emph{calx} means lime or chalk. English \emph{recalcitrant}, kicking back, comes from this one.

\begin{cousinweb}
\textbf{calx} \textquotedblleft{}heel\textquotedblright{} is a native Latin noun with no settled origin; Lewis \& Short keeps it apart from \emph{calx} \textquotedblleft{}lime\textquotedblright{}. English \emph{recalcitrant} comes from it, through \emph{recalcitrō} \textquotedblleft{}to kick back\textquotedblright{}.
\end{cousinweb}

\begin{grammarlens}[title={What you've built}]
One more word for telling, explaining and describing at B1.
\end{grammarlens}

\subsection*{Your turn}
\begin{itemize}
\raggedright
  \item \emph{Say it:} \emph{calx, calcis}
  \item \emph{Say it:} \emph{calx, calcis}, once more
  \item \emph{Say it:} a sentence of your own with \emph{calx, calcis}
  \item \emph{Recall:} say the Latin for side, then the Latin for hip, then say \emph{calx, calcis} again
\end{itemize}

\begin{tcolorbox}[breakable,colback=teal!4,colframe=teal!35!black,title={Before you move on}]
What does calx, calcis mean? (\textquotedblleft{}Heel\textquotedblright{}.) Say it once more.
\end{tcolorbox}
\section[planta, plantae]{planta, plantae — sole of the foot; a sprout}

\label{lesson:LA-C518-planta}



\begin{quote}
Before the new one: say the Latin for hip, then the Latin for heel.
\end{quote}

\subsection*{You'll want to know: planta, plantae}
\textbf{planta, plantae} — \textquotedblleft{}sole of the foot; a sprout\textquotedblright{}.

\begin{quote}
Puer nūdīs plantīs in harēnā ambulat.
\end{quote}

\begin{quote}
\emph{The boy walks on the sand with bare soles.}
\end{quote}

Feminine, first declension. Both the sole of the foot and a young plant. English \emph{plant}.

\begin{cousinweb}
\textbf{planta} is inherited from a Proto-Indo-European word for \textquotedblleft{}sole of the foot\textquotedblright{}, built on a root meaning \textquotedblleft{}flat\textquotedblright{} (EtymDB). English \emph{plant} comes from its second sense, a sprout.
\end{cousinweb}

\begin{grammarlens}[title={What you've built}]
One more word for telling, explaining and describing at B1.
\end{grammarlens}

\subsection*{Your turn}
\begin{itemize}
\raggedright
  \item \emph{Say it:} \emph{planta, plantae}
  \item \emph{Say it:} \emph{planta, plantae}, once more
  \item \emph{Say it:} a sentence of your own with \emph{planta, plantae}
  \item \emph{Recall:} say the Latin for hip, then the Latin for heel, then say \emph{planta, plantae} again
\end{itemize}

\begin{tcolorbox}[breakable,colback=teal!4,colframe=teal!35!black,title={Before you move on}]
What does planta, plantae mean? (\textquotedblleft{}Sole of the foot; a sprout\textquotedblright{}.) Say it once more.
\end{tcolorbox}
\section[umbilīcus, umbilīcī]{umbilīcus, umbilīcī — navel; the knob of a book-roll}

\label{lesson:LA-C518-umbilicus}



\begin{quote}
Before the new one: say the Latin for heel, then the Latin for sole of the foot.
\end{quote}

\subsection*{You'll want to know: umbilīcus, umbilīcī}
\textbf{umbilīcus, umbilīcī} — \textquotedblleft{}navel; the knob of a book-roll\textquotedblright{}.

\begin{quote}
Īnfāns parvum umbilīcum habet.
\end{quote}

\begin{quote}
\emph{The baby has a little navel.}
\end{quote}

Masculine, second declension. Also the knob on the rod of a book-roll, so \emph{ad umbilīcum addūcere} — to finish a book. English \emph{umbilical}.

\begin{cousinweb}
\textbf{umbilīcus} is inherited from a Proto-Indo-European word for \textquotedblleft{}navel\textquotedblright{}, and Greek \emph{omphalos} is a cousin (EtymDB; Lewis \& Short). English \emph{umbilical} comes from it, and \emph{navel} is a distant Germanic cousin.
\end{cousinweb}

\begin{grammarlens}[title={What you've built}]
One more word for telling, explaining and describing at B1.
\end{grammarlens}

\subsection*{Your turn}
\begin{itemize}
\raggedright
  \item \emph{Say it:} \emph{umbilīcus, umbilīcī}
  \item \emph{Say it:} \emph{umbilīcus, umbilīcī}, once more
  \item \emph{Say it:} a sentence of your own with \emph{umbilīcus, umbilīcī}
  \item \emph{Recall:} say the Latin for heel, then the Latin for sole of the foot, then say \emph{umbilīcus, umbilīcī} again
\end{itemize}

\begin{tcolorbox}[breakable,colback=teal!4,colframe=teal!35!black,title={Before you move on}]
What does umbilīcus, umbilīcī mean? (\textquotedblleft{}Navel; the knob of a book-roll\textquotedblright{}.) Say it once more.
\end{tcolorbox}
